\chapter{The strong-potential limit}
\label{chap:strong}

Along a line of potentials $v=\lambda u$ the Hamiltonian $A+\diag V+\lambda\diag w$, with
$w(y)=\sum_iu_in_i(y)$, is a matrix of polynomials in $\lambda$. Let $d:\Conf\to\N$ satisfy $d(c)=0$ and
$d(y)\le d(z)+1$ whenever $A_{yz}\neq0$, $y\ne z$ (a ``distance from $c$'').

\section{Interaction independence}

\begin{definition}
  \label{def:Hl}
  \lean{DFT.Hl, DFT.Lip, DFT.Hl_apply}
  \leanok
  $H_\lambda=A+\diag(V+\lambda w)$ as a polynomial matrix; $\mathrm{Lip}(A,d)$ is the growth condition above.
\end{definition}

\begin{theorem}[Degree bound and leading coefficient]
  \label{thm:entry_top}
  \lean{DFT.pow_succ_coeff, DFT.entry_bound, DFT.entry_top, DFT.entry_natDegree, DFT.entry_zero}
  \leanok
  \uses{def:Hl}
  $(H_\lambda^k)_{yc}$ has degree at most $k-d(y)$, and vanishes if $k<d(y)$. Its coefficient of degree
  $k-d(y)$ does not depend on $V$.
\end{theorem}
\begin{proof}
  \leanok
  Induction on $k$. An interaction factor always replaces a factor $\lambda$, so it only contributes below
  the top degree.
\end{proof}

\begin{theorem}[Interaction independence of the Krylov determinant]
  \label{thm:det_top}
  \lean{DFT.prod_top, DFT.det_top}
  \leanok
  \uses{thm:entry_top, def:kry}
  The coefficient of degree $\sum_j\mathrm{ord}(j)-\sum_yd(y)$ of $\det\Kry(H_\lambda,c)$ does not depend on $V$.
\end{theorem}
\begin{proof}
  \leanok
  Leibniz expansion; each term is a product of entries whose top coefficients are $V$-independent.
\end{proof}

\begin{theorem}[Reduction to free fermions]
  \label{thm:krylovPoly_ne_zero_of_top}
  \lean{DFT.Hl_eval, DFT.det_kry_eval, DFT.krylovPoly_ne_zero_of_top}
  \leanok
  \uses{thm:det_top, lem:eval_krylovPoly}
  If that coefficient is nonzero for $V=0$, then $P_c\neq0$ for every interaction $V$.
\end{theorem}
\begin{proof}
  \leanok
  A nonzero coefficient means a nonzero polynomial in $\lambda$, which has a non-root $\lambda_0$; evaluate
  $P_c$ at $\lambda_0u$.
\end{proof}

\section{Power sums}

\begin{definition}[Leading coefficients]
  \label{def:T}
  \lean{DFT.T, DFT.T_zero, DFT.T_succ}
  \leanok
  \uses{thm:entry_top}
  $T_k(y)$ is the coefficient of degree $k-d(y)$ of $(H_\lambda^k)_{yc}$ at $V=0$. It obeys
  $T_0(y)=[y=c]$ and $T_{k+1}(y)=w_yT_k(y)+\sum_{z:\,d(z)+1=d(y)}A_{yz}T_k(z)$.
\end{definition}

\begin{theorem}[Divided-difference representation]
  \label{thm:power_sum}
  \lean{DFT.power_sum, DFT.alphaRep, DFT.alphaRep_supp, DFT.alphaRep_rep}
  \leanok
  \uses{def:T}
  If the energies $w$ are pairwise distinct, then $T_k(y)=\sum_z\alpha_y(z)w_z^k$ for all $k$, with
  $\alpha_y$ supported on $\{y\}\cup\{z: d(z)<d(y)\}$.
\end{theorem}
\begin{proof}
  \leanok
  Strong induction on $d(y)$. The recursion prescribes $\alpha_y(z)$ for $z\ne y$, and $\alpha_y(y)$ is
  fixed by $T_0$.
\end{proof}

\begin{theorem}[Factorisation]
  \label{thm:top_coeff_factor}
  \lean{DFT.det_triangular_d, DFT.prod_coeff_top, DFT.top_coeff_factor, DFT.vandermonde_factor_ne_zero}
  \leanok
  \uses{thm:power_sum, thm:det_top}
  The leading coefficient of the free Krylov determinant equals $\big(\prod_y\alpha_y(y)\big)\det[w_z^{\mathrm{ord}(j)}]$.
  The Vandermonde factor is nonzero for distinct energies.
\end{theorem}
\begin{proof}
  \leanok
  $N=\alpha\cdot\mathrm{Vand}$, and $\alpha$ is triangular with respect to $d$. A permutation with all
  factors nonzero would strictly increase $\sum d$, so only the identity contributes.
\end{proof}

\begin{theorem}[(H$^*$) from the diagonal coefficients]
  \label{thm:hstar_of_alpha}
  \lean{DFT.hstar_of_alpha}
  \leanok
  \uses{thm:top_coeff_factor, thm:krylovPoly_ne_zero_of_top}
  If along some $u$ the energies are distinct and all $\alpha_y(y)\ne0$, then $P_c\neq0$ for every interaction.
\end{theorem}
\begin{proof}
  \leanok
  Combine the factorisation with Theorem~\ref{thm:krylovPoly_ne_zero_of_top}.
\end{proof}

\begin{theorem}[Uniqueness of power-sum coefficients]
  \label{thm:ps_unique}
  \lean{DFT.eval_weighted_sum, DFT.ps_unique}
  \leanok
  If $\sum_ia_ix_i^k=\sum_jb_jy_j^k$ for all $k$, then for every value $v$,
  $\sum_{x_i=v}a_i=\sum_{y_j=v}b_j$. Repeated nodes are allowed.
\end{theorem}
\begin{proof}
  \leanok
  Apply both sides to the Lagrange basis polynomial that equals $1$ at $v$ and $0$ at all other nodes.
\end{proof}

\section{Explicit coefficients}

\begin{theorem}[Recursion for the coefficients]
  \label{thm:alpha_rec}
  \lean{DFT.Ssrc, DFT.alpha_rec, DFT.alpha_sum, DFT.Ssrc_self}
  \leanok
  \uses{thm:power_sum, thm:ps_unique}
  $\alpha_y(z)(w_z-w_y)=\sum_{p:\,d(p)+1=d(y)}A_{yp}\alpha_p(z)$ and $\sum_z\alpha_y(z)=[y=c]$.
\end{theorem}
\begin{proof}
  \leanok
  Compare power-sum coefficients in $T_{k+1}=w_yT_k+\sum_pA_{yp}T_k(p)$.
\end{proof}

\begin{theorem}[Explicit diagonal coefficient]
  \label{thm:alpha_diag}
  \lean{DFT.Fres, DFT.alpha_diag, DFT.Fres_rec}
  \leanok
  \uses{thm:alpha_rec}
  With the resolvent route sum $F_y(x)=\sum_z\alpha_y(z)/(x-w_z)$:
  $\alpha_y(y)=[y=c]+\sum_pA_{yp}F_p(w_y)$, and $(x-w_y)F_y(x)=[y=c]+\sum_pA_{yp}F_p(x)$ away from poles.
\end{theorem}
\begin{proof}
  \leanok
  Partial fractions applied to the recursion of Theorem~\ref{thm:alpha_rec}.
\end{proof}
